\chapter{The Applications}\label{r04:ch:apps}

Two applications build RVV variants on the legacy stack: waLBerla's own
\code{UniformGridCPU} benchmark and a new benchmark app in waLBerla-wind. Both name the
instruction set per CMake target, so one build tree holds the scalar reference, upstream
\isa{rvv} and every \isa{rvv_ext} variant, compiled alike. This chapter describes both,
two bugs that only showed on the K1, and the Python environment.

\section{UniformGridCPU}\label{r04:sec:ugc}

\file{apps/benchmarks/UniformGridCPU} is upstream's lattice Boltzmann throughput
benchmark of the legacy core: a periodic box with a shear flow, all six streaming
patterns, several collision models. Its generator script reads
\code{<stencil>_<pattern>_<collision>} from \code{ctx.config} and asks for no explicit
vectors (\code{cpu_vec = {"instruction_set": None}}). Two additions, nothing removed:
\begin{description}[style=nextline]
  \item[\file{UniformGridCPU.py}, 3 lines] an optional fourth part of the config is the
    instruction set of the LBM kernels:
\begin{pycode}
    # Optional fourth part: vector instruction set of the LBM kernels, e.g. rvv or rvv_ext-vf-fma
    if len(config_tokens) > 3:
        cpu_vec = {"instruction_set": '_'.join(config_tokens[3:])}
\end{pycode}
    The \code{join} keeps names that contain \code{_}, such as \isa{rvv_ext}, whole.
  \item[\file{CMakeLists.txt}, 22 lines appended] a cache list
    \code{UNIFORMGRIDCPU_EXTRA_CONFIGS} of such configs, one generated library and one
    executable each, e.g.\ \code{UniformGridCPU_d3q19_aa_srt_rvv_ext-vf-fma}. The
    reference stays upstream's own target \code{UniformGridCPU_d3q19_pull_srt}.
  \item[\file{UniformGridCPU.cpp}] after the time loop, a checksum of the final state
    (\cref{r04:sec:checksum}).
\end{description}
Upstream adds this app only with \code{WALBERLA_BUILD_WITH_PYTHON=ON}
(\file{apps/benchmarks/CMakeLists.txt}), because it links \code{python_coupling}.
That option downloads pybind11 and compiles everything with \code{-fPIC}; all variants
of a tree are compiled alike, so their comparison stays fair. The app is run from a
\file{.prm} parameter file (\file{scripts-legacy/uniform_grid_cpu.prm}) instead of
upstream's Python setup scripts.

\section{The waLBerla-wind benchmark app}\label{r04:sec:wind-app}

waLBerla-wind's own applications are full wind-farm simulations with turbines, output
and snapshots. For throughput measurements a new app,
\file{apps/BenchmarkApp_cpu_uniform}, was made with waLBerla-wind's AppGenerator, as
its README intends for custom setups: CPU, uniform grid, with turbine output, VTK
output, snapshots and controllers switched off, so that the time loop holds the flow
solver only. Then:
\begin{itemize}
  \item \file{LBMGeneration.py} takes the variant \code{<pattern>_<instruction set>} from
    \code{ctx.config} (\code{pull_none}, \code{aa_rvv},
    \code{aa_rvv_ext-vf-fma-pf128}, \dots; \code{none} means no explicit vectors) and,
    for in-place patterns, generates the macroscopic getter as an even/odd pair, as the
    AppGenerator does for AA. The method is waLBerla-wind's: D3Q27 cumulant, Guo
    forcing, QR subgrid model, \code{fzyx}, float32, and \code{nontemporal} stores.
  \item \file{CMakeLists.txt} builds one executable per entry of the cache list
    \code{WIND_BENCHMARK_VARIANTS}.
  \item \file{main.cpp} runs \code{Parameters.warmupSteps} untimed steps before the
    timed ones and prints the checksum after the run.
  \item \file{benchmark.prm} has one $64^3$ block per process, periodic inflow and
    outflow and no turbine; shifted-periodic inflow refuses in-place patterns.
\end{itemize}
This gives waLBerla-wind AA streaming on the CPU without any change to waLBerla-wind's
library: the legacy core does the even/odd bookkeeping. On x86 pull and AA agree to
about $10^{-9}$ relative in float32 (\cref{r04:sec:app-checks}). Full AA support in
waLBerla-wind's production apps needs one more change that the benchmark does not:
its strain-rate writer reads the PDFs directly, which is wrong for the odd step.

\section{\texttt{OPTIMIZE\_FOR\_LOCALHOST} must be off}\label{r04:sec:localhost}

The first build of the benchmark app was configured with
\code{WALBERLA_OPTIMIZE_FOR_LOCALHOST=ON} and failed in code generation:
\begin{plaincode}
The expression cos(windDirection) of type cos is not known yet
\end{plaincode}
The option makes \code{config_from_context} give every kernel that names no
instruction set the best one of the build machine (\cref{r04:sec:where-isa}). That
includes waLBerla-wind's flow driver, whose kernel uses \code{cos} of a wind direction,
which \ps{}~1.4 cannot vectorize. On the K1 it would also have given every unnamed
kernel \isa{rvv}, so even the ``scalar'' variants would not have been scalar. So all
legacy builds here use \code{OPTIMIZE_FOR_LOCALHOST=OFF} and take the architecture
flags from the machine instead: \code{-march=rv64gcv_zvl256b} on the K1, what waLBerla's
CMake gives GCC there. (GCC~14.2 on the K1 rejects \code{-mtune=spacemit-x60}.)

\section{Bug 1: block-wise kernels and cache-line zeroing}\label{r04:sec:blockwise-bug}

The first K1 build of all nine waLBerla-wind variants stopped at \code{aa_rvv}:
\begin{plaincode}
waLBerlaWind_SweepCollection.cpp:5489: error: '_size_pdfs_0' was not declared in this scope
\end{plaincode}
(and the same for \code{_size_pdfs_1}, \code{_size_pdfs_2} and the three velocity
sizes). The cause is in waLBerla, not in our code, and needs three things together:
\begin{enumerate}
  \item An in-place pattern: only then does \code{lbmpy_walberla} generate the
    block-wise kernels (\cref{r04:ch:codegen}).
  \item Vector code with \code{nontemporal} and an instruction set that has
    \code{'cachelineZero'} (RVV, NEON, VSX). Then every written field's sizes are kernel
    parameters (\code{use_all_written_field_sizes}, \cref{r04:sec:nontemporal}); here
    \code{pdfs} and \code{velocity}.
  \item A kernel with more than one field. The template
    \file{python/lbmpy_walberla/templates/LBMSweepCollection.tmpl.cpp} declared the sizes
    of the \emph{first} field only (alphabetically \code{density}), instead of asking the
    kernel which it needs.
\end{enumerate}
\code{UniformGridCPU}'s block-wise kernel has only \code{pdfs}, so it was never
affected; waLBerla-wind's has four fields. The fix declares the sizes of every field of
the kernel. All fields of a block have the same interior size, and smaller sizes only
make the zeroing guard stricter, so the block's sizes are correct for all of them;
\code{[[maybe_unused]]} keeps kernels without the guard free of warnings:
\begin{r04diffcode}
    const int64_t _size_0 = size_0[level];
-   int64_t _size_{{block_stream_collide['all_fields'][0].name}}_0 = size_1;
-   int64_t _size_{{block_stream_collide['all_fields'][0].name}}_1 = size_2;
-   int64_t _size_{{block_stream_collide['all_fields'][0].name}}_2 = size_3;
+   {#- every field's sizes: the cacheline zeroing of vectorized nontemporal stores ... #}
+   {%- for field in block_stream_collide['all_fields'] %}
+   [[maybe_unused]] int64_t _size_{{field.name}}_0 = size_1;
+   [[maybe_unused]] int64_t _size_{{field.name}}_1 = size_2;
+   [[maybe_unused]] int64_t _size_{{field.name}}_2 = size_3;
+   {%- endfor %}
\end{r04diffcode}
Checked on the laptop before rebuilding on the board: \code{aa_rvv},
\code{aa_rvv_ext-vls256-vf-fma-pf128} and \code{aa_none} regenerated with the command
CMake uses and syntax-checked with the riscv64 cross compiler and the board's flags; no
errors and no new warnings, while the board's old file gives the 12 errors. On the
board all nine variants then built. Because the code generation reruns only when the
generator script changes (\cref{r04:sec:cmake}), the stale generated files of
\code{aa_rvv} had to be removed by hand.

\section{Bug 2: misaligned MPI buffers on the K1}\label{r04:sec:misaligned}

With the build fixed, every single-process run of every variant passed its checksum,
and every \emph{four}-process run died, the scalar \code{pull_none} included:
\begin{plaincode}
Signal: Bus error (7)
Signal code: Invalid address alignment (1)
\end{plaincode}
The stack trace points at the generated pack kernel that copies PDFs into the MPI send
buffer (\code{waLBerlaWind_StorageSpecification::PackKernels::packDirection}); the
faulting instruction is a vector segment store, \code{vsseg3e32.v}. Under gdb, on all
four ranks, the buffer addresses that waLBerla passes to \code{packDirection} are two
bytes off 4-byte alignment for every second block, and the store faults at such an
address with \code{vl} = 1. Three facts meet:
\begin{itemize}
  \item waLBerla's buffered communication writes a 6-byte header before each block's
    data (a length byte and, here, one byte of block ID, then the 4-byte
    \code{Direction} enum), safely with \code{memcpy}. The generated pack kernel then
    writes floats through \code{reinterpret_cast<float*>(outBuffer)} (upstream template
    \file{LbmStorageSpecification.tmpl.cpp}). C++ does not allow a misaligned
    \code{float*}; x86 tolerates it, and so do the X60's scalar stores (6.49\,ns against
    6.39\,ns per load and store in a test on one core).
  \item GCC auto-vectorizes this scalar \ps{} code when \code{-march} contains the V
    extension: the loops of the 12 edge directions (3 PDFs per cell) become
    \code{vsseg3e32.v}, those of the 8 corners \code{vse32.v}, while the 6 faces (9 PDFs
    per cell) stay scalar \code{fsw}. That is why the misaligned face packs pass and the
    first edge faults.
  \item The X60 requires element alignment for vector accesses: \code{vle32} at byte
    offset 2 and \code{vle64} at offset 4 raise a bus error, \code{vle32} at offset 4 or
    12 (element-aligned, not vector-aligned) works. The RVV specification allows either
    behaviour. Linux 6.6 on the board emulates only scalar misaligned accesses (the
    signal's address is the faulting instruction's, as its misaligned-trap path reports
    it) and does not know the \code{riscv_hwprobe} key for vector misaligned access.
\end{itemize}
A 30-line program reproduces it outside waLBerla: the same loop writing into a byte
buffer faults at offsets 2 and 6 with \code{-march=rv64gcv} (with or without
\code{zvl256b}), runs at offsets 0 and 4, and runs at every offset with
\code{-march=rv64gc}.

Neither our changes nor the board's setup create the defect: none of our waLBerla
commits touch the communication path, and waLBerla-wind's own applications use the same
pack info and scheme. What exposes it is compiling with the V extension, which every RVV
build needs. The board's GCC defaults to \code{rv64imafdc} without V, and upstream
waLBerla's \code{OPTIMIZE_FOR_LOCALHOST} adds \code{-march=native}, which GCC~14 rejects
on RISC-V; so an upstream build on the K1 contains no vector code and never meets the
bug. It affects every multi-process run of a legacy waLBerla application on the K1
built with the V extension, \code{UniformGridCPU} too.

A fix for upstream waLBerla is prepared in \file{docs/bug-fix-pr/walberla-packinfo-buffer-alignment/},
as a patch against the release v7.4 that also applies to master. Upstream had fixed the
same problem for the pack infos generated by \code{pystencils_walberla} in 2024 and
documented in \code{UniformPackInfo} that every implementation must align its data; the
patch adds \code{SendBuffer::forwardAligned} and \code{RecvBuffer::skipAligned}, which do
that with the same scheme, and uses them in the three pack infos that did not. Tested on
upstream's own test \code{InterpolationNoSlip} on two processes: 108 misaligned accesses
under the alignment sanitizer on x86 and a bus error on the K1 before the patch, none
after it. The folder also holds a script that reproduces both. Whether this project's
waLBerla branch takes the patch for the K1 runs was not decided when this was written.

\section{The Python environment}\label{r04:sec:env}

The legacy code generation needs its own virtual environment, \file{.venv-legacy}
(next to the workspace, covered by the root \file{.gitignore} rule \code{.venv*/}):
\begin{itemize}
  \item Python 3.12, SymPy 1.12.1 (\ps{}~1.4's upper limit), NumPy, jinja2 for the
    templates, py-cpuinfo for detection on x86, \code{lbmpy==1.4} installed without its
    dependencies (it would otherwise pull in a released \ps{}), pytest~8 (pytest~9 cannot
    load \ps{}~1.4's \file{conftest.py}).
  \item \ps{} itself is not installed: \code{PYTHONPATH=submodules/pystencils/src} at
    configure time makes the build use the submodule (\cref{r04:sec:cmake}).
  \item On the K1: PyPI has no NumPy for riscv64, so the venv is created with
    \code{--system-site-packages} and uses Ubuntu's NumPy 1.26.4. The root file system
    with \file{/tmp} is almost full, so compilers and pip get \code{TMPDIR} in the work
    folder.
\end{itemize}
\code{scripts-legacy/bench.py setup} creates this environment on any machine.
